% =====================================================================
%  Chapter 8  ·  Session 8 --- DevSecOps, WAF and secure applications
% =====================================================================
\chapter{DevSecOps, WAF and secure applications}

The machine, the network, the transport and the identities are secured, and
the application can still be the hole. Session~6 showed why: an
SQL-injection flaw hands an attacker the data through the front door, and no
encryption or firewall helps, because the application itself does the
leaking. This chapter therefore covers the application layer from two sides:
security built into the development process (DevSecOps), and a shield in
front of the running application (the web application firewall).

\section{Application security in the cloud}

Application security protects applications from threats across their whole
life, from design to production, and the cloud changes it in three ways
(CSA, 2024, p.~230). First, applications are built as groups of
microservices and external services communicating over APIs, so the attack
surface is spread across many small components and their interfaces rather
than one monolith. Secondly, they are developed with rapid DevOps practices,
which is a risk, since fast change can outrun review, and an opportunity,
since automation can enforce security on every change. Thirdly, they are
assembled from libraries written by others, much of it open source, which
brings in \emph{supply-chain} risk, because a vulnerability in a dependency
one did not write is still a vulnerability in one's application. The
countermeasure to the third is the \emph{software bill of materials} (SBOM),
an inventory of every component and dependency, produced by software
composition analysis in the image pipeline (CSA, 2024, pp.~181--182). When a
library is found vulnerable, the SBOM answers the one question that matters
during an incident: whether and where the library is in use.

\begin{realworld}{Log4Shell, 2021}
In December 2021 a critical flaw was disclosed in Log4j, a logging library so
common that it sat, often invisibly, inside a large share of the world's Java
applications. A single crafted string written to a log could make a vulnerable
system fetch and run attacker code: trivial to exploit and broad in reach. The
main difficulty was not the patch but finding out where Log4j was used at
all: most organisations did not know which of their applications and
third-party products embedded it. A software bill of materials answers that
question in advance. Knowing one's dependencies before the need arises is a
control in its own right.
\end{realworld}

\section{The secure development lifecycle}

A flaw caught in design costs a fraction of the same flaw caught in
production, so security work is moved earlier (CSA, 2024, pp.~231--232).

\begin{definition}[Secure development lifecycle]\label{def:ssdlc}
The \textbf{secure software development lifecycle} (SSDLC) describes the
integration of security into every stage of building software: secure
\emph{design and architecture}, secure \emph{build, integration and testing},
secure \emph{deployment}, and \emph{runtime defence and monitoring} of the
application in production. Moving security work earlier is called
\emph{shifting left}.
\end{definition}

Two practices give the lifecycle substance. The first is \emph{threat
modelling} at design time, a structured process to identify, assess and
remediate threats before the system is built (CSA, 2024, p.~232). The
Guidance recommends \emph{STRIDE}, which sorts threats into six categories,
Spoofing, Tampering, Repudiation, Information disclosure, Denial of service
and Elevation of privilege (CSA, 2024, pp.~232--233), and thereby turns the
open-ended threat question of Session~1 into a checklist that is hard to
leave gaps in. The second practice is automated security testing in two
forms. \emph{Static} application security testing (SAST) examines the source
code before deployment for security flaws and logic errors and scans for
hard-coded credentials, keys and tokens before they enter the repository
(CSA, 2024, p.~234); it produces false positives and therefore needs
tuning. \emph{Dynamic} testing (DAST) probes the running application from
outside afterwards, emulating real attacks (CSA, 2024, p.~235). Static
testing sees the code but not its behaviour, and dynamic testing sees the
behaviour but not the code, so the two are used together, as symmetric and
asymmetric encryption were in Session~6.

\section{DevSecOps}

The lifecycle becomes practical when it is automated into the pipeline that
already builds and ships the software.

\begin{definition}[CI/CD and DevSecOps]\label{def:devsecops}
A \textbf{continuous integration / continuous delivery} (CI/CD) pipeline
describes the set of automated steps that build, test and deploy software as
developers merge changes. \textbf{DevSecOps} (development, security and
operations) describes embedding security into that pipeline, so that automated
security checks run on every change rather than as a separate, occasional
gate.
\end{definition}

Comer (2021) describes the CI/CD machinery and two deployment strategies
that make frequent releases safe (Comer, 2021, pp.~211--213). \emph{Canary
deployment} sends the new version to a small fraction of users first and
observes before rolling on, and \emph{blue/green deployment} keeps two
environments so that a release can be switched over, and back, instantly.
DevSecOps adds security to each phase: SAST runs during continuous
integration, so that vulnerable code does not merge, and DAST runs during
continuous delivery against the staged application (CSA, 2024,
pp.~243--245). This closes a loop with Session~5, because the immutable
infrastructure of the cloud-native world is what makes the pipeline
trustworthy: every environment is built fresh from the same tested
definition, so there is no drift between what was tested and what runs. The
CSA framing, six pillars, is cultural as much as technical: security is a
shared responsibility across development, operations and security,
automated to minimise human error and measured for continuous improvement
(CSA, 2024, p.~244). In contrast to a security team that reviews at the end,
DevSecOps places the check where the change is made.

\begin{figure}[htbp]
\centering
\begin{tikzpicture}[font=\footnotesize]
  \def\R{2.55}
  \tikzset{stage/.style={draw=kgrey, fill=klight, rounded corners=2pt,
      inner sep=3pt, minimum width=1.35cm, align=center},
    arr/.style={-{Stealth[length=2mm]}, kristiania, line width=0.7pt}}
  % eight stages, clockwise from the top
  \node[stage] (plan)    at (90:\R)   {Plan};
  \node[stage] (code)    at (45:\R)   {Code};
  \node[stage] (build)   at (0:\R)    {Build};
  \node[stage] (test)    at (-45:\R)  {Test};
  \node[stage] (release) at (-90:\R)  {Release};
  \node[stage] (deploy)  at (-135:\R) {Deploy};
  \node[stage] (operate) at (180:\R)  {Operate};
  \node[stage] (monitor) at (135:\R)  {Monitor};
  % clockwise flow
  \draw[arr] (plan)    to[bend left=12] (code);
  \draw[arr] (code)    to[bend left=12] (build);
  \draw[arr] (build)   to[bend left=12] (test);
  \draw[arr] (test)    to[bend left=12] (release);
  \draw[arr] (release) to[bend left=12] (deploy);
  \draw[arr] (deploy)  to[bend left=12] (operate);
  \draw[arr] (operate) to[bend left=12] (monitor);
  \draw[arr] (monitor) to[bend left=12] (plan);
  % security at the centre
  \node[draw=kristiania, fill=kristiania!8, circle, align=center,
        font=\footnotesize\itshape, text=kristiania, minimum size=1.9cm]
        at (0,0) {Security\\at every\\stage};
  % where the two automated tests sit
  \node[font=\scriptsize, text=kgrey, anchor=west] at (0:\R) [xshift=7mm] {SAST};
  \node[font=\scriptsize, text=kgrey, anchor=west] at (-45:\R) [xshift=6mm] {DAST};
\end{tikzpicture}
\caption{The DevSecOps loop. Development (plan, code, build, test) flows
continuously into operations (release, deploy, operate, monitor) and back;
DevSecOps embeds automated security at every stage rather than as a gate at
the end, with SAST as the code is built and DAST against the deployed
application.}
\label{fig:devsecops}
\end{figure}

What lets a pipeline build the same tested definition every time is that
the infrastructure itself is written down as text. \emph{Infrastructure as
code} (IaC) replaces hand-configuration, which no two engineers perform
identically, with a machine-readable specification that provisions the
environment automatically (Comer, 2021, pp.~121--122). Comer connects it to
\emph{zero touch provisioning}, bringing systems up with no manual steps,
and separates an \emph{imperative} specification, which lists the steps to
perform, from a \emph{declarative} one, which states the desired end state
and lets the tooling reach it, the same work-towards-a-goal style as the
Kubernetes control plane of Session~5. Three security benefits follow.
Consistency removes the configuration drift that breeds vulnerabilities. The
specification is reviewable and version-controlled, so a change is visible,
auditable and revertible, which Session~12 uses for compliance. And because
the environment is rebuilt from the text rather than patched by hand, it is
immutable in the sense of Session~5. IaC and pipelines are, however,
themselves part of the attack surface and need the same access control as
the application (CSA, 2024, p.~238).

\section{The web application firewall}

Secure development reduces the flaws that reach production, but it does not
remove all of them, and it does nothing for flaws in third-party code one
does not control. The running application therefore needs a guard of its
own.

\begin{definition}[Web application firewall]\label{def:waf}
A \textbf{web application firewall} (WAF) describes a component placed in
front of a web application that inspects incoming HTTP requests at the
application layer and blocks those that match patterns of attack, such as SQL
injection, cross-site scripting and the application-layer floods of
Session~3, before they reach the application.
\end{definition}

Here the distinction of Session~4 applies: an NGFW secures the network
broadly, a WAF secures one web application specifically, and they operate at
different layers. The WAF is the control that answers the SQL-injection
example of Session~6, because it inspects a request for the shape of an
injection attempt and blocks it before the application can execute it. In
the SSDLC it belongs to the runtime-defence phase, the external shield
around a production application that complements secure code within. The
Guidance accordingly treats it as a control that prevents attacks from
arriving, not as a way to repair a flawed application or to excuse careless
development (CSA, 2024, p.~247). On the students' service it runs in front
of Nginx; in Azure it is the Application Gateway's WAF at Layer~7.

\begin{example}[A WAF for Nextcloud]\label{ex:waf-trc}
The SQL-injection scenario of Session~6 now meets a control that fits it. \emph{Threat}: a crafted request exploits an injection flaw, possibly in a
third-party Nextcloud plugin, to read files from the database, bypassing TLS
and the firewall because it arrives as an ordinary-looking HTTPS request.
\emph{Risk}: the likelihood is non-trivial, because plugins and dependencies
carry such flaws, and the impact is high, because the confidentiality of every
stored file is at stake. \emph{Controls}, in depth: secure coding and dependency scanning in the
pipeline reduce the chance that the flaw exists (shift left); an SBOM shows
at once whether a newly disclosed library vulnerability affects the service;
and a WAF in front of Nextcloud inspects and blocks the malicious request at
runtime even if the flaw is present. The layers are independent by design: the WAF catches what the
pipeline missed, and the pipeline shrinks what the WAF must catch.
\end{example}

\paragraph{Reading for this chapter.} CSA Security Guidance, Domain~10, in
particular \S10.1 (Secure Development Lifecycle, threat modelling and
STRIDE, pre- and post-deployment testing, pp.~231--236), \S10.4 (DevSecOps:
CI/CD \& Application Testing, pp.~243--248) and \S10.5 (Serverless \&
Containerized Application Considerations, pp.~248--250); the domain runs
pp.~230--252. Domain~8, \S\S8.1.4--8.1.5 (Software Composition Analysis and
the Software Bill of Materials, pp.~181--182). Comer, Chapter~15 (DevOps:
CI/CD, canary and blue/green deployment, pp.~207--214) and \S\S9.10--9.11
(infrastructure as code; declarative, imperative and intent-based
specifications, pp.~121--122). Optional deepening: J\o sang, Chapter~11
(Secure by Design, pp.~243--270). The Microsoft Learn module for Session~8 (Azure
Application Gateway and its WAF) is linked on Canvas.

\section{Summary}

The application is a layer no network control reaches, and the SQL
injection of Session~6 shows why it needs its own defences. In the cloud
those defences must cope with microservices and APIs, rapid change and
third-party dependencies, the last tracked with a software bill of
materials. Security is built in through the secure development lifecycle
(Definition~\ref{def:ssdlc}), which models threats with STRIDE and tests
with SAST before deployment and DAST after, and it is made routine by
embedding those checks in the CI/CD pipeline as DevSecOps
(Definition~\ref{def:devsecops}), which the immutable, code-defined
infrastructure of Session~5 keeps trustworthy. Because flaws still reach
production and third-party code cannot be fixed at will, a web application
firewall (Definition~\ref{def:waf}) shields the running application and
blocks the injection attack at runtime (Example~\ref{ex:waf-trc}). Every
control so far can still fail, and Session~9 therefore turns to detection.

\section*{Self-check}
\addcontentsline{toc}{section}{Self-check}

Sketch answers follow the exercises.

\begin{exercise}[Shifting left]
What does shifting left mean, and why is a vulnerability cheaper to fix in
design than in production?
\end{exercise}

\begin{exercise}[SAST vs DAST]
Distinguish SAST from DAST. Why are they used together, and which earlier
pair of complementary techniques does this resemble?
\end{exercise}

\begin{exercise}[What DevSecOps adds]
Explain what DevSecOps adds to a CI/CD pipeline, and how the immutable
infrastructure of Session~5 makes the pipeline's testing trustworthy.
\end{exercise}

\begin{exercise}[NGFW vs WAF, revisited]
An NGFW and a WAF both inspect traffic. State the difference in one sentence,
and name the one that blocks an SQL-injection attempt against a Nextcloud
login.
\end{exercise}

\begin{exercise}[A disclosed library flaw]
A newly disclosed vulnerability affects a popular open-source library. Which
document shows immediately whether the application is affected, and which
runtime control could block exploitation even if it is?
\end{exercise}

\paragraph{Sketch answers.}
{\small
(1)~Shifting left moves security work earlier in the lifecycle (design and
build rather than production); a flaw caught in design is fixed by changing a
plan, while the same flaw in production may mean an incident, an emergency
patch and lost data. (2)~SAST examines source code before deployment to find
flaws; DAST probes the running application from outside afterwards. Each has a
blind spot, SAST seeing code but not behaviour and DAST behaviour but not
code, so they are combined, as symmetric and asymmetric encryption are in
Session~6. (3)~DevSecOps embeds automated security checks (SAST in CI, DAST in
CD) into every change rather than an occasional gate; because immutable
infrastructure builds every environment fresh from one tested definition,
there is no drift between what was tested and what runs. (4)~An NGFW secures
the network broadly with application awareness; a WAF secures one web
application by inspecting HTTP requests; the WAF blocks the SQL-injection
attempt. (5)~A software bill of materials (SBOM) shows whether and where the
library is used; a web application firewall could block exploitation of the
flaw at runtime even if the vulnerable library is present.}


